\section{The integral Koszul invariant behind the key rule}
\label{sec:ftb-koszul-key}
This section translates the algebraic construction used by
\srcid{ExistsIntegralKeyInvariantRuleCharZero}.
The construction is independent of a compactification.
The input is a principally polarized abelian scheme \(A/R\) of
relative dimension \(g\geq1\), a symmetric ample principal line bundle
\(L\), a generator of \(L|_0\), a fibrewise nonzero theta section
\(\vartheta\in H^0(A,L)\), and a basis \(v_1,\ldots,v_g\) of its Lie
algebra.  This is the source datum \(\operatorname{KeyIn}(R)\).
In particular the principal polarization has trivial translation
kernel.  Changing the theta section or Lie basis does not change the
rigidification.

We use the following abelian-scheme laws with their indicated scope.
For \(m>0\), \(H^0(A,L^m)\) is free of rank \(m^g\) over a local base
and its displayed bases commute with base change; positive powers have
the higher-cohomology vanishing used in their affine Cech complexes.
For \(m\geq4\) these sections separate points and first-order tangent
vectors on a field fibre.  Multiplication by \(m\) is finite locally
free of degree \(m^{2g}\), is etale if \(m\) is invertible, and has
\(m^{2g}\) points over an algebraically closed field with \(m\ne0\).
We also use the theorem of the cube, symmetric multiplication
\([m]^*L\simeq L^{m^2}\), and faithfully flat descent along \([m]\).
These are precisely the laws used below; a scheme equipped only with
an ample line bundle would not be sufficient.

\subsection{Squared torsion of the actual augmented complex}
Let \(C=\bigoplus_{j\geq0}C_j\) be a finite free graded module over a
commutative ring, with fixed finite bases, and let \(D\) have degree
\(-1\), \(D^2=0\).
A contraction is a degree \(+1\) matrix \(h\) with \(Dh+hD=1\).
When the even and odd basis sets are identified by a bijection \(e\),
put
\[
 \tau(D,h,e)
  =\det\bigl((D+h)_{\mathrm{even},e(\mathrm{even})}\bigr)^2.
 \tag{K1}\label{eq:ftb-torsion-square}
\]
The square is part of the definition and removes the ordering sign.

\begin{srlemma}[Integral contraction and unique squared torsion]
\label{lem:ftb-squared-torsion}
Over a local ring \(R\), if the residue-field complex is exact,
a contraction exists, the even and odd ranks are equal, and
\eqref{eq:ftb-torsion-square} is a unit of \(R\).
It is independent of \(h\) and \(e\).
\end{srlemma}
\begin{proof}
Over a field choose \(G\) with \(DGD=D\), by splitting
\(\operatorname{im}D\).  Exactness gives
\(N^2=0\) for \(N=1-(DG+GD)\): \(DN=ND=0\), and the image of
\(N\), being in \(\ker D=\operatorname{im}D\), is killed by \(N\).
Thus \(DG+GD\) is invertible.
Lift \(G\) entry by entry from the residue field to \(H\) over \(R\).
The matrix \(B=DH+HD\) has unit determinant and commutes with \(D\).
Then \(h_0=HB^{-1}\) satisfies \(Dh_0+h_0D=1\).
Keep only the degree \(+1\) entries of \(h_0\).  The degree-zero
part of the last identity is still \(Dh+hD=1\), and the other
entries vanish for degree reasons.

The matrix \(D+hDh\) is an involution and exchanges even and odd
degrees: \(DhD=D\), and
\(Dh(hD)=0\).  Its two off-diagonal blocks are inverse, proving
equality of the two finite ranks.
For independence take two contractions \(h_1,h_2\) and write
\(k=h_2-h_1\).  Then
\[
 Dk+kD=0,\quad (kD)^2=0,\quad
 (D+h_2)(D+h_1)=1+kD+h_2h_1 .
\]
Order the even basis by degree.  This last matrix is block triangular;
on degree \(j\) its diagonal block is \((1+kD)|_{C_j}\).
Every such block has determinant one.  For \(j=0\), the block \(kD\)
is zero.  For successive degrees use
\(\det(1+XY)=\det(1+YX)\), with
\(X=k:C_{j-1}\to C_j\) and \(Y=D:C_j\to C_{j-1}\);
the swapped product is \(-kD|_{C_{j-1}}\).
Since that product squares to zero, the determinants of
\(1+kD\) and \(1-kD\) are inverse.  Induction gives determinant one
on every degree.

Fix the same odd ordering for both contractions.
If \(a,b\) are the determinants of their even-to-odd blocks, and
\(c\) the determinant of the odd-to-even block for \(h_1\), the
previous paragraph gives \(ac=bc=1\).  Hence \(a=b\) and \(a\)
is a unit.  A different odd ordering changes a determinant only by
the sign of a permutation.  Squaring proves all assertions.
\end{proof}

Let \(s_1,\ldots,s_g\in H^0(A,L^n)\).
For an integer \(l>gn\), choose bases \(b_j\) of
\(H^0(A,L^{l-jn})\), \(0\leq j\leq g\).
The source complex has
\[
 C_{j+1}=\bigoplus_{|I|=j}H^0(A,L^{l-jn}),
 \qquad C_0=R^d,
 \tag{K2}\label{eq:ftb-koszul-complex}
\]
with the usual signed multiplication by the \(s_i\) and the final
evaluation map at \(z_1,\ldots,z_d\).
In increasing-index exterior coordinates, deletion of \(i\) has
sign \((-1)^{|\{a\in I:a<i\}|}\).
This fixes the convention used in the determinant calculations.

The source transverse-zero hypothesis means all of the following:
the \(z_a\) are sections through pairwise distinct points in
\emph{every} fibre; their images are exactly the common zeros of
the \(s_i\); a line-bundle chart is chosen at every \(z_a\);
\(s_i(z_a)=0\); and
\[
 J_a=\det\bigl(D_{v_k}s_i(z_a)\bigr)_{i,k}
                 \in R^\times .                   \tag{K3}\label{eq:ftb-jac-unit}
\]
Values and derivatives here are the coefficients in those chosen
line-bundle charts.

\begin{srlemma}[The transverse Koszul unit]\label{lem:ftb-koszul-unit}
Let \(R\) be local Noetherian and let the displayed transverse-zero
hypothesis hold.  Then the complex
\eqref{eq:ftb-koszul-complex} is exact on every algebraically closed
field fibre.  It has a unique squared torsion \(t_l\in R^\times\),
and
\[
                     d=g!\,n^g.                   \tag{K4}\label{eq:ftb-zero-count}
\]
\end{srlemma}
\begin{proof}
On a field fibre, the Jacobian condition says that the germs of the
\(s_i\) generate the maximal ideal at each listed zero.
At other points the zero ideal is the unit ideal.
The local Koszul complex is therefore exact.
Here is the source passage to global sections: refine to a finite
affine cover trivializing \(L\), use its finite-intersection Cech
cube, and successively cut by the \(s_i\).
Weak regularity on each affine intersection gives the local
division step.  The Cech complexes of the powers
\(L^{l-jn}\) have no positive cohomology because each exponent is
positive.  Induction through these cuts shows (i) exactness at every
positive Koszul degree, (ii) interpolation of arbitrary residues at
the finite zeros, and (iii) that a section vanishing at those zeros
is \(\sum_i s_i a_i\) with \(a_i\in H^0(A,L^{l-n})\).
For (ii), first interpolate the prescribed residues separately on
affine charts; their differences on overlaps lie in the ideal
\((s_1,\ldots,s_g)\), and the preceding Cech division steps glue
them.  For (iii), the same division steps glue local expressions
in that ideal.  These three assertions give exactness of the
augmented evaluation complex, including its degree-zero term.
This is the source Cech/Koszul argument, with the exponents
\(l-jn>0\) retained.

All matrices and evaluation charts commute with field base change.
Thus exactness over the algebraic closure of the residue field
descends to the residue field (apply a linear retraction to a
solution of each finite linear system).
Lemma~\ref{lem:ftb-squared-torsion} supplies \(t_l\).
Equality of the even and odd ranks gives
\[
 d=\sum_{j=0}^g(-1)^j\binom gj(l-jn)^g=g!\,n^g.
\]
The last identity is the \(g\)-th finite difference of \(X^g\);
lower powers cancel and the leading power contributes \(g!n^g\).
\end{proof}

\subsection{The split torsion determinant}
For each positive integer \(m\) choose an integral matrix
\(\alpha_m\) of order four satisfying
\(\alpha_m^t\alpha_m=mI_4\).
It can be chosen for every \(m\): write \(m=a^2+b^2+c^2+d^2\)
and use
\[
 \alpha_m=
 \begin{pmatrix}
 a&-b&-c&-d\\ b&a&-d&c\\
 c&d&a&-b\\d&-c&b&a
 \end{pmatrix}.
 \tag{K5}\label{eq:ftb-quaternion}
\]
The relation also gives \(\alpha_m\alpha_m^t=mI_4\).

Suppose \(R\) is a local Noetherian principal ideal domain, \(m\) is
a unit, and a fibrewise distinct full \(N\)-torsion enumeration
\(\xi_1,\ldots,\xi_{N^{2g}}\) is supplied, with \(m\mid N\).
The kernel of \(\alpha_m^t:A^4\to A^4\) has
\(m^{4g}\) elements in a geometric fibre.
One sees the number by Smith normal form:
if \(d_i\) are its four invariant factors, then \(d_i\mid m\),
\(\prod_i d_i=|\det\alpha_m|=m^2\), and the kernel has
\(\prod_i d_i^{2g}=m^{4g}\) points.
Choose an enumeration
\[
 \zeta_I=(\zeta_{I1},\ldots,\zeta_{I4}),
 \qquad I\in\{1,\ldots,m^g\}^{\,4}.
\]
These are actual tuples of the supplied sections.
Indeed every such kernel tuple is \(m\)-torsion, since
\(\alpha_m\alpha_m^t=mI\), and hence is \(N\)-torsion.
Reduction maps the supplied \(N\)-torsion bijectively to the
geometric \(N\)-torsion.  A relation true after reduction is true
over \(R\): two sections of the finite etale \(N\)-torsion scheme
that agree at the closed point agree on the connected local base.
Thus
\[
 \sum_r\alpha_{m,rc}\zeta_{Ir}=0\quad(c=1,\ldots,4)
 \tag{K6}\label{eq:ftb-kernel-tuples}
\]
and distinct \(I\)'s stay distinct in every fibre.

Let \(a_{Ir}:A^4\to A\) be
\(a_{Ir}(x)=\sum_c\alpha_{m,rc}x_c+\zeta_{Ir}\).
The theorem of the cube gives a generator \(w_I\) of
\[
 \left(\bigotimes_r a_{Ir}^*L\right)
       \otimes\left(\bigotimes_c\operatorname{pr}_c^*L^{-m}\right).
 \tag{K7}\label{eq:ftb-box-generator}
\]
For precision, in the Picard group the quadratic law for a symmetric
line bundle cancels all off-diagonal terms by
\(\alpha_m^t\alpha_m=mI\).
The linear cross terms in the translations cancel by
\eqref{eq:ftb-kernel-tuples}.  Constant line bundles pulled from
\(\operatorname{Spec}R\) are free because \(R\) is local.
These are exactly the terms in the source cube calculation.

Normalize a generator
\(\mu:[m]^*L\otimes L^{-m^2}\simeq\mathcal O_A\)
at the origin using the chosen rigidification.
Since each \(\zeta_{Ir}\) is \(m\)-torsion, the product of its
values and the origin value of \(w_I\) has the same line-bundle
weight as the rigidification to exponent \(4-4m^3\).
Their quotient is therefore a unit \(\rho_I\in R^\times\), namely
\[
 w_I(0)^{m^2}\prod_r\mu(\zeta_{Ir})
              =\rho_I\,\operatorname{rig}^{\,4-4m^3}.
 \tag{K8}\label{eq:ftb-rho-normalization}
\]
This identity is an equality of framed values, so is independent of
the temporary chart used to write it.

For a basis \(b_1,\ldots,b_{m^g}\) of \(H^0(A,L^m)\), divide
\(\prod_r a_{Ir}^*\vartheta\) by \(w_I\) and expand it in the
external-product basis.  Denote the coefficient matrix by \(Q_m\):
\[
 \prod_r a_{Ir}^*\vartheta
   =w_I\sum_{I'}(Q_m)_{II'}\prod_c\operatorname{pr}_c^*b_{I'_c}.
 \tag{K9}\label{eq:ftb-box-matrix}
\]
All its coefficients are in \(R\), by the relative basis and
base-change law.  The source box-product spanning proof obtains this
basis by affine base change and then one factor at a time.

\begin{srlemma}[Unit determinant, with its actual kernel argument]
\label{lem:ftb-box-det-unit}
The determinant of \(Q_m\) is a unit.  Consequently
\[
        \lambda_m=(\det Q_m)^{2m^2}\prod_I\rho_I^2
                         \in R^\times              \tag{K10}\label{eq:ftb-lambda}
\]
is defined integrally.
\end{srlemma}
\begin{proof}
It suffices to reduce to an algebraic closure \(k\) of the residue
field.  Put \(H_m=\ker(\alpha_m:A^4(k)\to A^4(k))\), and choose
one \(I_0\).  The ratios \(w_{I_0}/w_I\) transform under translation
by \(H_m\) by characters \(\chi_I\).  The numerator theta products
in \eqref{eq:ftb-box-matrix} are invariant under \(H_m\), since
all their arguments are unchanged.  Thus those products multiplied
by \(w_{I_0}/w_I\) are eigenvectors of characters \(\chi_I\).

The characters are pairwise distinct.  If two agreed, pull their
invariant ratio back along the \(r\)-th row embedding
\(A\to A^4\), \(x\mapsto(\alpha_{m,r1}x,\ldots,\alpha_{m,r4}x)\).
The relation \(\alpha_m\alpha_m^t=mI\) cancels the other rows.
After a constant trivialization this gives an isomorphism
\[
 [m]^*\bigl(t_{\zeta_{Ir}}^*L\otimes
                     t_{\zeta_{I'r}}^*L^{-1}\bigr)\simeq\mathcal O_A
\]
invariant under \(A[m]\).
Faithfully flat descent along \([m]\) gives
\(t_{\zeta_{Ir}}^*L\simeq t_{\zeta_{I'r}}^*L\).
The principal-polarization kernel is trivial, so
\(\zeta_{Ir}=\zeta_{I'r}\) for every row, hence \(I=I'\).

Each eigenvector is nonzero.  Choose a point where \(\vartheta\)
is nonzero.  The isogeny \(\alpha_m\) is surjective (its composite
with \(\alpha_m^t\) is \([m]\)); one can therefore prescribe four
nonzero theta arguments after the translation.
Eigenvectors with distinct multiplicative characters are linearly
independent over a field: in a shortest nontrivial relation,
translate by a group element distinguishing two characters and
subtract an appropriate multiple of the relation to shorten it.
A nonzero left-kernel vector of \(Q_m\), combined with
\eqref{eq:ftb-box-matrix} and multiplied by \(w_{I_0}/w_I\), would
contradict this independence.  Thus \(\det Q_m\ne0\) over \(k\).
A coefficient over the local ring with nonzero residue is a unit.
Equation~\eqref{eq:ftb-lambda} is now a product of units.
\end{proof}

\subsection{The exact formula and cancellation of the auxiliary choices}
Choose \(n,l,l'\) with \(n\geq4\), \(gn<l<l'\), and write
\[
 D=\prod_{j=0}^g(l-jn)(l'-jn),\quad s=3g+2,\quad d=g!n^g.
\]
Supply split \(D\)-torsion and the transverse tuple above, and assume
\(D\in R^\times\).  Apply
Lemma~\ref{lem:ftb-box-det-unit} to every degree \(l-jn,l'-jn\).
Let \(t_l,t_{l'}\) be the two squared Koszul units,
\(\lambda_j,\lambda'_j\) the corresponding units
\eqref{eq:ftb-lambda}, and \(J=\prod_{a=1}^dJ_a\).
Set
\[
 \begin{split}
 T_l&=t_l^{\,4D^s}
       \prod_{j=0}^g\lambda_j^{
          (-1)^j\binom gj(D/(l-jn))^s},\\
 T_{l'}&=t_{l'}^{\,4D^s}
       \prod_{j=0}^g(\lambda'_j)^{
          (-1)^j\binom gj(D/(l'-jn))^s},\\
 \kappa&=T_l^{\,gn-2l'}T_{l'}^{\,2l-gn}
                         J^{\,8D^s(l'-l)},\\
 c&=2D^s\,d(l'-l).
 \end{split}                                      \tag{K11}\label{eq:ftb-kaform}
\]
All negative powers in this formula are powers of units.

\begin{srproposition}[Integral value and its frame law]
\label{prop:ftb-integral-koszul-invariant}
The formula \eqref{eq:ftb-kaform} gives a unit of \(R\), and
commutes with local base change preserving the displayed data.
Over a characteristic-zero field, it is independent of the bases,
enumeration and charts of the fixed transverse intersection, and of
the auxiliary generators used in \eqref{eq:ftb-box-generator}.
If \(\vartheta\) is replaced by \(a\vartheta\) and the Lie basis by
the matrix \(M\), its value satisfies
\[
              a^{8c}\kappa_{\rm new}
                         =(\det M)^{4c}\kappa_{\rm old}.
 \tag{K12}\label{eq:ftb-koszul-homog}
\]
Here the tuple \(s_i\) is fixed; independence of \emph{that} tuple
is a further step, stated separately below.
\end{srproposition}
\begin{proof}
Existence over the stated local principal ideal domain follows from
the preceding lemmas, which also give \eqref{eq:ftb-zero-count}.
Base change applies to the actual multiplication and evaluation
matrices, their contractions, the cubical generators, and determinants;
the defining identities therefore base change term by term.

We record the determinant calculations, including their signs.
If a basis in degree \(m\) changes by a matrix of determinant \(b\),
its fourth external product has determinant \(b^{4m^{3g}}\).
Thus
\[
             \lambda_{\rm new}
                   =\lambda_{\rm old}b^{-8m^{3g+2}}.
\]
Rescaling a cubical generator by a unit \(z_I\) rescales row \(I\)
of \(Q_m\) by \(z_I^{-1}\), and rescales \(\rho_I\) by
\(z_I^{m^2}\); these cancel in \eqref{eq:ftb-lambda}.
A permutation of kernel tuples contributes only a determinant sign,
removed by the square.  For arbitrary two choices of kernel enumeration, first permute the
full kernel tuples to agree.  The normalized generators \(\mu\)
agree: their quotient is a base unit, determined by its origin value.
On the proper geometrically connected variety \(A^4\), corresponding
\(w_I\)'s have a constant unit ratio \(z_I\).
The origin normalization gives
\(\rho'_I=z_I^{m^2}\rho_I\); uniqueness of the external-product
coefficients gives the inverse row factor \(z_I^{-1}\) in \(Q_m\).
These factors cancel as above.  External-product independence follows
one factor at a time from independence of the original basis, so this
comparison is valid without choosing an analytic realization.

For a change of bases \(b_j\) with determinants \(b_j\), the squared
torsion changes by
\(\prod_jb_j^{2(-1)^j\binom gj}\).
If the chosen line generator at each zero changes by \(g_a\), set
\(\gamma=\prod_ag_a\), in the convention that old evaluations are
\(g_a^l\) times new evaluations in degree \(l\).
The augmented degree-zero evaluation term then gives
\[
 t_{l,\rm new}
   =t_{l,\rm old}\gamma^{-2l}
                    \prod_jb_j^{2(-1)^j\binom gj},
 \qquad J_{\rm new}=J_{\rm old}\gamma^{-gn}.
\]
The torsion comparison follows by conjugating the whole graded
matrix, with these evaluation scalars in degree zero, and taking
the determinant of the even-to-odd block.  Permuting the zero
sections again disappears after squaring.

In \(T_l\) the basis exponents cancel term by term:
\[
 4D^s\cdot2(-1)^j\binom gj
 -(D/m)^s(-1)^j\binom gj\cdot8m^s=0.
\]
Hence \(T_{l,\rm new}=T_l\gamma^{-8lD^s}\), and the analogous
formula holds for \(l'\).
The remaining exponent of \(\gamma\) in \(\kappa\) is
\[
 -8D^s\{l(gn-2l')+l'(2l-gn)+gn(l'-l)\}=0 .
\]
This proves the claimed independence for a fixed intersection.

Multiplication of \(\vartheta\) by \(a\) multiplies every row of
\(Q_m\) by \(a^4\).  There are \(m^{4g}\) rows; thus
\(\lambda_m\) is multiplied by \(a^{8m^{4g+2}}\).
The same finite-difference identity as in
\eqref{eq:ftb-zero-count} gives
\(T_l\mapsto a^{8D^sd}T_l\), and similarly for \(l'\).
The total exponent is \(-16D^sd(l'-l)=-8c\).
Changing the Lie basis multiplies each \(J_a\) by \(\det M\);
there are \(d\) zeros, giving exponent
\(8D^sd(l'-l)=4c\).
These are exactly \eqref{eq:ftb-koszul-homog}.
\end{proof}

\begin{srlemma}[Numerical data away from a residue prime]
\label{lem:ftb-key-numerical-data}
For every \(g\geq1\) and every prime \(p\), choose
\[
 n=4p,\qquad l=gn+1,\qquad l'=gn+1+p
\]
and the matrices \eqref{eq:ftb-quaternion}.
Then \(n\geq4\), \(gn<l<l'\), \(D\geq1\), \(c\geq1\),
and \(p\nmid D\).
\end{srlemma}
\begin{proof}
For \(0\leq j\leq g\),
\[
 l-jn=4p(g-j)+1,\qquad
 l'-jn=p(4(g-j)+1)+1 .
\]
Both are positive and congruent to one modulo \(p\).
Their product \(D\) is therefore positive and prime to \(p\);
the formula for \(c\) is a product of positive integers.
\end{proof}

\paragraph{What this construction supplies and what remains.}
The proposition supplies the actual integral \(\operatorname{KaVal}\)
from split torsion and split transverse zeros, including the exact
normalizations and powers.  It does not yet assert a rule independent
of the chosen sections.  The source next proves this independence by
DVR specialization along lines, first between algebraically independent
tuples and then between arbitrary admissible tuples; it then proves
holomorphy in period coordinates, normalizes the powered rule, and
descends its exponent to one by connected families.
Those are separate producer steps.  In particular neither the unit
determinant in \eqref{eq:ftb-lambda} nor the unit
\eqref{eq:ftb-kaform} alone is the exponent-one key unit used in the
carrying section.
